\chapter{Matriz de operacionalizaci\'on de variables}

\begin{table}[H]
\centering
\caption{Matriz de operacionalizaci\'on de variables del estudio}
\label{tab:anexo_operacionalizacion}
\footnotesize
\begin{tabularx}{\textwidth}{p{2cm}p{2.8cm}p{1.8cm}p{2.5cm}p{1.5cm}}
\toprule
\textbf{Variable} & \textbf{Definici\'on operacional} & \textbf{Dimensi\'on} & \textbf{Indicador} & \textbf{Escala} \\
\midrule
Precio FOB & Valor unitario promedio de exportaci\'on de palta Hass & Econ\'omica & USD/kg mensual & Continua \\
Volumen exportado & Cantidad de palta Hass exportada por periodo & Comercial & Toneladas/mes & Continua \\
Mes / temporada & Periodo calendario y ventana comercial peruana & Temporal & Mes (1--12) & Ordinal \\
Rezago precio & Valores hist\'oricos rezagados de la serie FOB & Temporal & USD/kg ($t-1$, $t-2$) & Continua \\
Destino principal & Pa\'is de destino con mayor participaci\'on & Comercial & C\'odigo pa\'is & Nominal \\
Utilidad percibida & Valoraci\'on del productor sobre la herramienta & Cognitiva & Puntaje Likert 1--5 & Ordinal \\
Precisi\'on predictiva & Error del modelo seleccionado & T\'ecnica & RMSE, MAPE & Continua \\
Margen proyectado & Diferencia entre ingresos y costos simulados & Econ\'omica & Soles/ha & Continua \\
\bottomrule
\end{tabularx}
\parbox{0.94\textwidth}{\small \textit{Nota.} Elaboraci\'on propia a partir de las variables del modelo predictivo (Cap\'itulo I) y los indicadores de evaluaci\'on (Cap\'itulo II).}
\end{table}
